\begin{titlepage}
  \begin{center}

  {\Huge Mini Colecovision}

  \vspace{25mm}

  \includegraphics[width=0.90\textwidth,height=\textheight,keepaspectratio]{img/SPARKLETRON.png}

  \vspace{25mm}

  \today

  \vspace{15mm}

  {\Large Jay Convertino}

  \end{center}
\end{titlepage}

\tableofcontents

\newpage

\section{Introduction}

\par
This project creates a mini colecovision for a portable console. Much of the TTL and Analog Monostable circuits
are emulated by a CPLD. The full PCB and CPLD code is in this repository. It emulates the original Colecovision
with the additional super game module.

\subsection{Specifications}

\par
\begin{itemize}
  \item Z80 CPU
  \item 32 KiB of RAM
  \item 32 kiB of ROM
  \item SN76489 Sound Chip
  \item YMZ284 Sound Chip
  \item TMS9118 Video Display Processor with 16 KiB of VRAM
\end{itemize}

\subsection{Parts List}
\begin{longtable}{ |*{4}{c|} }
\hline
{Item} & {Qty} & {Reference(s)} & {Value} \\
\hline
{1} & {13} & {C1, C7, C8, C10 to C14, C18, C19, C23, C25, C29} & {100nF} \\
\hline
{2} & {1} & {C2} & {330uF} \\
\hline
{3} & {4} & {C3, C5, C24, C30} & {10uF} \\
\hline
{4} & {1} & {C4} & {100pF} \\
\hline
{5} & {1} & {C6} & {270pF} \\
\hline
{6} & {8} & {C9, C17, C22, C31 to C35} & {100nF} \\
\hline
{7} & {2} & {C15, C16} & {33pF} \\
\hline
{8} & {2} & {C20, C21} & {10nF} \\
\hline
{9} & {1} & {C26} & {0.47uF} \\
\hline
{10} & {1} & {C27} & {0.1uF} \\
\hline
{11} & {1} & {C28} & {2.2uF} \\
\hline
{12} & {1} & {D1} & {LED} \\
\hline
{13} & {5} & {H1, H2, H4, H6, H7} & {MountingHole\_Pad} \\
\hline
{14} & {1} & {J1} & {Conn\_01x07} \\
\hline
{15} & {3} & {J2, J10, J11} & {Conn\_01x02} \\
\hline
{16} & {1} & {J3} & {Conn\_Coaxial} \\
\hline
{17} & {1} & {J4} & {Cartridge Port} \\
\hline
{18} & {1} & {J5} & {Player 1} \\
\hline
{19} & {1} & {J6} & {Conn\_02x05\_Odd\_Even} \\
\hline
{20} & {1} & {J7} & {Player 2} \\
\hline
{21} & {1} & {J9} & {SJ1-3525NG} \\
\hline
{22} & {3} & {L1, L2, L3} & {4.7uH} \\
\hline
{23} & {1} & {Q1} & {2N3904} \\
\hline
{24} & {1} & {R1} & {4k7} \\
\hline
{25} & {1} & {R2} & {470R} \\
\hline
{26} & {2} & {R3, R28} & {100k} \\
\hline
{27} & {2} & {R4, R27} & {100K} \\
\hline
{28} & {2} & {R5, R6} & {2k2} \\
\hline
{29} & {3} & {R7, R8, R9} & {3K3} \\
\hline
{30} & {1} & {R10} & {75R} \\
\hline
{31} & {1} & {R11} & {510R} \\
\hline
{32} & {1} & {R12} & {100R} \\
\hline
{33} & {1} & {R13} & {3k3} \\
\hline
{34} & {4} & {R14, R15, R16, R17} & {1k} \\
\hline
{35} & {2} & {R18, R19} & {10k} \\
\hline
{36} & {1} & {R20} & {1K} \\
\hline
{37} & {1} & {R21} & {1k} \\
\hline
{38} & {1} & {R22} & {220R} \\
\hline
{39} & {1} & {R23} & {10K} \\
\hline
{40} & {2} & {R24, R26} & {1K} \\
\hline
{41} & {1} & {R25} & {68K} \\
\hline
{42} & {2} & {RN1, RN2} & {10k} \\
\hline
{43} & {1} & {RV1} & {10K} \\
\hline
{45} & {1} & {SW1} & {SW\_Push} \\
\hline
{46} & {1} & {SW2} & {SW\_SPDT} \\
\hline
{47} & {1} & {SW3} & {SW\_SPDT} \\
\hline
{48} & {1} & {U1} & {TPA711D} \\
\hline
{49} & {1} & {U2} & {SN76489AN} \\
\hline
{50} & {1} & {U3} & {Z84C0010AEG} \\
\hline
{51} & {1} & {U4} & {CY62256-55PC} \\
\hline
{52} & {1} & {U5} & {27C256} \\
\hline
{53} & {1} & {U6} & {TMS9118NL} \\
\hline
{54} & {2} & {U7, U8} & {TMS4416} \\
\hline
{55} & {1} & {U9} & {EPM7128SLC} \\
\hline
{56} & {1} & {U10} & {74ABT125} \\
\hline
{57} & {1} & {U11} & {YMZ284} \\
\hline
{58} & {5} & {U12, U13, U14, U15, U16} & {74AHCT1G08} \\
\hline
{59} & {1} & {Y1} & {10.738635 MHz} \\
\hline
\end{longtable}

\section{Building}

\par
Steps to build the mini Colecovision

\begin{itemize}
  \item Create main PCB schematic/gerber/coleco\_original.zip
  \item Create Right Angle PCB schematic/gerber/right\_angle/right\_angle.zip
  \item Program ROM with BIOS
  \item Populate main PCB
  \item Populate right angle PCB
  \item Power up and program CPLD
  \item Build your own case
\end{itemize}
\par

\subsection{Dependecies}

\par
The following are the dependecies needed to build the firmware and PCB for the system.

\begin{itemize}
  \item Quartus 13.0 sp1
  \item python 3.X
  \item KiCAD v7.X
\end{itemize}

\subsubsection{Protable Coleco Glue File List}
\begin{itemize}
\item src
	\begin{itemize}
	\item {'src/porta\_glue\_coleco.v': {'file\_type': 'verilogSource'}}
	\end{itemize}
\item constr
	\begin{itemize}
	\item {'constr/porta\_glue\_coleco.sdc': {'file\_type': 'SDC'}}
	\end{itemize}
\item tb
	\begin{itemize}
	\item {'tb/tb\_porta\_glue\_coleco.v': {'file\_type': 'verilogSource'}}
	\end{itemize}
\end{itemize}


\subsubsection{Fusesoc}
\subsubsection{Protable Coleco Glue Targets}
\begin{itemize}
\item default
	\begin{itemize}
	\item[$\space$] Info: Default IP target for future tool intergration.
	\item src
	\item constr
	\end{itemize}
\item sim
	\begin{itemize}
	\item[$\space$] Info: Simulation target for basic test bench.
	\item src
	\item tb
	\end{itemize}
\end{itemize}


\par
Fusesoc is used for the simulation target only. There are no build targets due to the use of Quartus 13.0sp1

\subsubsection{Quartus}
\par
This project uses the last version that supports the MAX7000S series. The version is 13.0sp1 for quartus.
The project is located at src/quartus13sp01/.

\subsection{PCB}

\par
The four layer PCB is fairly easy to populate. The right angle PCB is a dual layer PCB which is even easier.
I recommend starting with resistors, then IC's, and then the rest. Surface mount parts should be done last.

\subsection{3D Printed Case}

\par
A 3D printed case model is not included. I've kept this for release in the future.

\subsection{Programming}

\par
There are two devices that need to be programmed.

\subsubsection{ROM}

\par
A minipro TL866A+ is an execelent device for programming the ROM with a BIOS.

\subsubsection{CPLD}

\par
Quartus 13.0sp1 is the easiest way to build and program the MAX7000 CPLD. You will need an altera blaster.
I recommend the chinese clone blasters, they actually worked the best. While the worst was the Terasic blaster
which did not work at all.

\newpage

\section{Usage}

\subsection{Directory Guide}

\par
Below highlights important folders from the root of mini\_colecovision.

\begin{enumerate}
  \item \textbf{docs} Contains all documentation related to this project.
    \begin{itemize}
      \item \textbf{datasheets} Contains all datasheets for components.
      \item \textbf{manual} Contains user manual and github page that are generated from the latex sources.
    \end{itemize}
  \item \textbf{img} Contains images of the project
  \item \textbf{schematic} KiCAD v7.X schematic and PCB designs
    \begin{itemize}
      \item \textbf{gerber} Contains gerber files and archives for production.
      \item \textbf{pdf} PDF schematic
    \end{itemize}
  \item \textbf{src} CPLD firmware source
    \begin{itemize}
      \item \textbf{protable\_coleco} Contains verilog source code and constraits
      \item \textbf{quartus13sp01} Quartus project to use to generate firmware file.
    \end{itemize}
\end{enumerate}

\section{Module Documentation}
